\documentclass{article}

\usepackage{amsmath,amssymb}
\usepackage{xurl}
\usepackage[hidelinks]{hyperref}
\setlength{\emergencystretch}{2em}

% Shared commands for notes in docs/paper-gaps/.

\DeclareUnicodeCharacter{2080}{\ifmmode{}_{0}\else\textsubscript{0}\fi}
\DeclareUnicodeCharacter{2081}{\ifmmode{}_{1}\else\textsubscript{1}\fi}
\DeclareUnicodeCharacter{2082}{\ifmmode{}_{2}\else\textsubscript{2}\fi}
\DeclareUnicodeCharacter{2099}{\ifmmode{}_{n}\else\textsubscript{n}\fi}

\newcommand{\Lean}{\textsc{Lean}}
\ExplSyntaxOn
\NewDocumentCommand{\leanid}{m}{
  \ifmmode
    \textsf{\detokenize{#1}}
  \else
    \group_begin:
    \urlstyle{sf}
    \tl_set:Nn \l_tmpa_tl {#1}
    \tl_replace_all:Nnn \l_tmpa_tl {\_} {_}
    \exp_args:NV \path \l_tmpa_tl
    \group_end:
  \fi
}
\ExplSyntaxOff
\providecommand{\tr}{\operatorname{tr}}
\newcommand{\tnleanIssueBase}{https://github.com/LionSR/TNLean/issues/}
\newcommand{\tnleanPrBase}{https://github.com/LionSR/TNLean/pull/}
\newcommand{\tnleanDocBase}{https://sirui-lu.com/TNLean/docs/}
\newcommand{\ghissue}[1]{\href{\tnleanIssueBase#1}{GitHub issue \##1}}
\newcommand{\ghpr}[1]{\href{\tnleanPrBase#1}{PR \##1}}
\newcommand{\doclink}[1]{\href{\tnleanDocBase#1.html}{\path{#1}}}

% Dirac notation and the identity operator, as in blueprint/src/macros/common.tex.
% \providecommand keeps note-local definitions authoritative.
\usepackage{dsfont}
\providecommand{\ket}[1]{|#1\rangle}
\providecommand{\bra}[1]{\langle #1|}
\providecommand{\braket}[2]{\langle #1 | #2 \rangle}
\providecommand{\ketbra}[2]{|#1\rangle\!\langle #2|}
\providecommand{\Id}{\mathds{1}}
\providecommand{\one}{\mathds{1}}
\providecommand{\C}{\mathbb{C}}
\providecommand{\MN}[1]{M_{#1}(\C)}
\providecommand{\MD}{M_D(\C)}

% External referencing for paper sources.  The build helper writes sanitized
% auxiliary files under build/paper-aux/ and excludes duplicated source labels.
\usepackage{xr}

% Import a generated paper auxiliary file with a caller-supplied label prefix.
% The candidate paths support builds from docs/paper-gaps/, the repository root,
% and build/paper-aux/ itself.
\newcommand{\importpaperaux}[2]{%
  \IfFileExists{../../build/paper-aux/#2.aux}{%
    \externaldocument[#1]{../../build/paper-aux/#2}%
  }{%
    \IfFileExists{build/paper-aux/#2.aux}{%
      \externaldocument[#1]{build/paper-aux/#2}%
    }{%
      \IfFileExists{#2.aux}{%
        \externaldocument[#1]{#2}%
      }{}%
    }%
  }%
}

% General external reference macros
\newcommand{\paperref}[2]{\ref{#1-#2}}
\newcommand{\papereqref}[2]{(\ref{#1-#2})}

% CPSV16 source (1606.00608 / MPDO-22-12-17-2.tex) external document and shortcuts
\importpaperaux{cpsv16-}{1606.00608/MPDO-22-12-17-2-tnlean}
\newcommand{\cpsvref}[1]{\paperref{cpsv16}{#1}}
\newcommand{\cpsveqref}[1]{\papereqref{cpsv16}{#1}}

% Verdict marker: \gapnote{<kind>}{<status>}. Kind is the policy
% classification (clarification, local-correction, scope-restriction,
% unfaithful, false-source, open-gap); status is open, wip (elimination
% underway), resolved, or historical. Machine-read by the site index;
% prints a verdict line.
\newcommand{\gapnote}[2]{\par\noindent\textbf{Verdict:} #1 (#2).\par}


\title{Block-Diagonal Boundary Conditions in the Degenerate Parent-Hamiltonian Ground Space}
\author{}
\date{2026-06-11}

\begin{document}
\maketitle
\gapnote{open-gap}{resolved}

\section{Notation}

Let $A=(A_i)_{i=1}^d$ be a tensor of bond dimension $D$. For $N\geq 1$ and
$X\in M_D(\mathbb C)$, define
\begin{align}
  \Gamma_N^A(X)(i_1,\ldots,i_N)
    &=
    \tr(A_{i_1}\cdots A_{i_N}X).
  \label{eq:cpgsv21_block_diagonal_parent_ground_space_gamma}
\end{align}
The associated open-boundary space and periodic matrix product vector are
\begin{align}
  \mathcal G_N^A
    &=
    \{\Gamma_N^A(X):X\in M_D(\mathbb C)\},
  \label{eq:cpgsv21_block_diagonal_parent_ground_space_open_space}\\
  V^{(N)}(A)
    &=
    \Gamma_N^A(\Id).
  \label{eq:cpgsv21_block_diagonal_parent_ground_space_periodic_vector}
\end{align}

Consider a ring of length $N$ with periodic boundary conditions. Let
$\operatorname{res}_{s,L}$ denote restriction to the length-$L$ interval
beginning at $s\in\mathbb Z/N\mathbb Z$. The corresponding local ground-space
condition is
\begin{align}
  \mathcal K_{N,L}(A)
    &=
    \bigcap_{s\in\mathbb Z/N\mathbb Z}
    \{\psi:\operatorname{res}_{s,L}(\psi)\in\mathcal G_L^A\}.
  \label{eq:cpgsv21_block_diagonal_parent_ground_space_chain_space}
\end{align}

For blocks $A^1,\ldots,A^b$ and nonzero weights
$\mu_1,\ldots,\mu_b$, define the block-diagonal tensor by
\begin{align}
  \widehat A_i
    &=
    \bigoplus_{j=1}^b\mu_jA_i^j,
  \label{eq:cpgsv21_block_diagonal_parent_ground_space_weighted_tensor}\\
  \mu_j
    &\neq 0
    \qquad (1\leq j\leq b).
  \label{eq:cpgsv21_block_diagonal_parent_ground_space_nonzero_weights}
\end{align}
The letters $g$ and $b$ both denote the number of blocks below: $g$ follows
the block-injective canonical-form notation of
Ref.~\cite{Cirac2021Matrix}, whereas $b$ follows Theorem~12 of
Ref.~\cite{PerezGarcia2007MPSReps}.

\section{The source assertion}

Section~IV.C of Ref.~\cite{Cirac2021Matrix} describes the degenerate
parent-Hamiltonian ground space of a matrix product state with more than one
block in canonical form.\footnote{For traceability, the boundary-crossing
comparison has a formal span-dependent form on the main branch
(\path{TNLean/MPS/ParentHamiltonian/BNTBlockDiagonalCrossingTrace.lean} and
\path{TNLean/MPS/ParentHamiltonian/BNTBlockDiagonalPGVWCComparison.lean}).
The short crossing-tail span hypothesis is removed by the global change-of-cut
comparison in
\leanid{MPSTensor.exists_blockDiagonal_boundary_chainGroundSpace_of_global_cut_bnt_c1_pgvwc07_of_dualFixedPoint};
this resolves
\href{https://github.com/LionSR/TNLean/issues/3597}{issue \#3597} at the
length bound of the source. The removal of the former extra $(L_0+1)$ in the
separation constant resolves
\href{https://github.com/LionSR/TNLean/issues/4195}{issue \#4195}. The earlier
tracker
\href{https://github.com/LionSR/TNLean/issues/2971}{issue \#2971}, which
tracked the boundary-crossing comparison itself, is now closed. Pull
request~\#2724 records the earlier conditional reduction on the main branch,
which reduced the unconditional target to boundary-crossing comparison
identities. Earlier reductions were tracked by
\href{https://github.com/LionSR/TNLean/issues/2641}{issue \#2641},
\href{https://github.com/LionSR/TNLean/issues/905}{issue \#905}, and
\href{https://github.com/LionSR/TNLean/issues/870}{issue \#870}. The relevant
source lines are
\path{Papers/2011.12127/TN-Review-main.tex}:2112--2128.}

The block-injective canonical form in Ref.~\cite{Cirac2021Matrix} gives a
block-diagonal physical--virtual correspondence. Write
\begin{align}
  \widetilde A^i
    &=
    \bigoplus_{j=1}^g A_j^i.
  \label{eq:cpgsv21_block_diagonal_parent_ground_space_canonical_tensor}
\end{align}
For every block-diagonal virtual matrix $X$, the correspondence states
\begin{align}
  X=(X_j)_j
    &\in
    \bigoplus_{j=1}^g M_{D_j}(\mathbb C),
  \label{eq:cpgsv21_block_diagonal_parent_ground_space_virtual_matrix}\\
  X
    &=
    \sum_i c_i(X)\widetilde A^i
    \quad\text{for some coefficients }(c_i(X))_i.
  \label{eq:cpgsv21_block_diagonal_parent_ground_space_physical_virtual}
\end{align}
After blocking, this correspondence becomes the eventual word-span identity
\begin{align}
  \operatorname{span}
  \{(A_1^\omega,\ldots,A_g^\omega):|\omega|=m\}
    &=
    \prod_{j=1}^g M_{D_j}(\mathbb C)
    \qquad (m\geq m_0)
  \label{eq:cpgsv21_block_diagonal_parent_ground_space_eventual_tuple_span}
\end{align}
for some $m_0$.

The formal development separates the product-span input from the finite
blocking hypothesis. It first proves the one-step intersection identity from
Theorem~12 of Ref.~\cite{PerezGarcia2007MPSReps} under
Eq.~\ref{eq:cpgsv21_block_diagonal_parent_ground_space_eventual_tuple_span}.
The normalized basis-of-normal-tensors (BNT) theorems then derive this product
span from injectivity of
\begin{align}
  \Gamma_{L_0}^{A_j}
    &:M_{D_j}(\mathbb C)
      \longrightarrow
      (\mathbb C^d)^{\otimes L_0}
  \label{eq:cpgsv21_block_diagonal_parent_ground_space_block_injectivity}
\end{align}
and the separated normalized-block hypotheses. For $b\geq 2$, the simultaneous
products span the product algebra at every length
\begin{align}
  n
    &\geq
    3(b-1)(L_0+1).
  \label{eq:cpgsv21_block_diagonal_parent_ground_space_tuple_span_range}
\end{align}
Thus the boundary-decomposition route does not replace Condition C1 of
Ref.~\cite{PerezGarcia2007MPSReps} by a length-one span condition. Instead, it
compares boundary conditions after a global change of cut and uses only the
simultaneous product span on the complementary segment of length $N-L_0$.

The parent-Hamiltonian assertion in Section~IV.C of
Ref.~\cite{Cirac2021Matrix} is
\begin{align}
  N\geq L_0+1
    \quad\Longrightarrow\quad
  \ker(H_N)
    &=
    \operatorname{span}
    \{V^{(N)}(A_j):j=1,\ldots,g\}.
  \label{eq:cpgsv21_block_diagonal_parent_ground_space_source_kernel}
\end{align}
The sentence preceding this theorem restricts the proof to block-diagonal
boundary conditions. After opening the periodic boundary at a chosen cut,
these boundary matrices have the form
\begin{align}
  X
    &=
    \bigoplus_{j=1}^g X_j,
  \label{eq:cpgsv21_block_diagonal_parent_ground_space_boundary_x}\\
  Y
    &=
    \bigoplus_{j=1}^g Y_j,
  \label{eq:cpgsv21_block_diagonal_parent_ground_space_boundary_y}\\
  X_j,Y_j
    &\in
    M_{D_j}(\mathbb C).
  \label{eq:cpgsv21_block_diagonal_parent_ground_space_boundary_blocks}
\end{align}

\section{Normalization scope}

The normalization used in the block-intersection calculation of
Ref.~\cite{PerezGarcia2007MPSReps} is
\begin{align}
  \sum_a A^j_a(A^j_a)^\dagger
    &=
    \Id.
  \label{eq:cpgsv21_block_diagonal_parent_ground_space_source_unital}
\end{align}
For the dual transfer operator, the general canonical form has a positive
definite fixed point $\Lambda_j$ satisfying
\begin{align}
  \sum_a (A^j_a)^\dagger\Lambda_j A^j_a
    &=
    \Lambda_j.
  \label{eq:cpgsv21_block_diagonal_parent_ground_space_dual_fixed_point}
\end{align}
The unrestricted source-length declarations assume
Eqs.~\ref{eq:cpgsv21_block_diagonal_parent_ground_space_source_unital}
and~\ref{eq:cpgsv21_block_diagonal_parent_ground_space_dual_fixed_point},
together with positivity of the dual fixed-point data. Unlike the diagonal
coordinate choice in Ref.~\cite{PerezGarcia2007MPSReps}, their proofs do not
require $\Lambda_j$ to be diagonal.

A temporary square-root gauge makes each block left-canonical only within the
fixed-length simultaneous product-span proof. Gauge invariance then transports
the span to the original source representatives. Long-word propagation, the
global change-of-cut comparison, the block-diagonal boundary calculation, and
the parent-Hamiltonian kernel argument all remain in the original
source-unital gauge. In particular, the formalization does not claim that a
nonunitary gauge preserves both normalization identities.

The declarations ending in
\texttt{\_pgvwc07\_of\_dualFixedPoint} prove the unrestricted chain-space,
boundary-closing, kernel-containment, and exact kernel-equality conclusions of
Theorem~12 in Ref.~\cite{PerezGarcia2007MPSReps}. The older declarations ending
in \texttt{\_pgvwc07} remain available for the doubly normalized specialization
$\Lambda_j=\Id$. Their main chain-space capstone follows by applying the
unrestricted theorem to the identity dual fixed point. This resolves the
normalization restriction tracked by
\href{https://github.com/LionSR/TNLean/issues/5751}{issue \#5751}, with the
source-range tuple-span layer supplied by
\href{https://github.com/LionSR/TNLean/issues/5764}{issue \#5764} and the
global-cut and kernel threading supplied by
\href{https://github.com/LionSR/TNLean/issues/5770}{issue \#5770}.

\section{The block-diagonal intersection identity}

The theorem entitled \emph{Degeneracy of the ground space v2} in
Ref.~\cite{PerezGarcia2007MPSReps} gives the finite-chain identity used by the
source argument.\footnote{Source location:
\path{Papers/quant-ph_0608197/MPSarchive.tex}:1424--1456.}
It identifies the block counts and imposes the local-length bound
\begin{align}
  b
    &=
    g,
  \label{eq:cpgsv21_block_diagonal_parent_ground_space_block_count}\\
  L
    &\geq
    3(b-1)(L_0+1)+1.
  \label{eq:cpgsv21_block_diagonal_parent_ground_space_local_length}
\end{align}
The intersection step applies in the range
\begin{align}
  m
    &\geq
    L.
  \label{eq:cpgsv21_block_diagonal_parent_ground_space_intersection_range}
\end{align}
In that range, the asserted identity is
\begin{align}
  \mathbb C^d\otimes
  \left(\bigoplus_j\mathcal G_m^{A_j}\right)
  \cap
  \left(\bigoplus_j\mathcal G_m^{A_j}\right)\otimes\mathbb C^d
    &=
    \bigoplus_j\mathcal G_{m+1}^{A_j}.
  \label{eq:cpgsv21_block_diagonal_parent_ground_space_intersection_identity}
\end{align}

Let $\psi$ belong to the left-hand side of
Eq.~\ref{eq:cpgsv21_block_diagonal_parent_ground_space_intersection_identity}.
The two decompositions of $\psi$ are
\begin{align}
  \psi
    &=
    \sum_j\sum_{i_1,\ldots,i_{m+1}}
    \tr\!\left(A^j_{i_{m+1}}C^j_{i_1}
      A^j_{i_2}\cdots A^j_{i_m}\right)
    \ket{i_1\cdots i_{m+1}},
  \label{eq:cpgsv21_block_diagonal_parent_ground_space_left_decomposition}\\
  \psi
    &=
    \sum_j\sum_{i_1,\ldots,i_{m+1}}
    \tr\!\left(D^j_{i_{m+1}}A^j_{i_1}
      A^j_{i_2}\cdots A^j_{i_m}\right)
    \ket{i_1\cdots i_{m+1}}.
  \label{eq:cpgsv21_block_diagonal_parent_ground_space_right_decomposition}
\end{align}
The direct-sum separation lemma and one-block injectivity then give, for every
$j$, $i_1$, and $i_{m+1}$,
\begin{align}
  A^j_{i_{m+1}}C^j_{i_1}
    &=
    D^j_{i_{m+1}}A^j_{i_1}.
  \label{eq:cpgsv21_block_diagonal_parent_ground_space_cd_identity}
\end{align}

The corresponding source line in Ref.~\cite{PerezGarcia2007MPSReps} writes
$E^j=\sum_k C^j_kA^j_k$ without an adjoint on the second factor. Because the
next line uses
Eq.~\ref{eq:cpgsv21_block_diagonal_parent_ground_space_source_unital}, the
calculation requires the adjoint-corrected definition
\begin{align}
  E^j
    &=
    \sum_k C^j_kA^{j\dagger}_k.
  \label{eq:cpgsv21_block_diagonal_parent_ground_space_corrected_e}
\end{align}
Using Eq.~\ref{eq:cpgsv21_block_diagonal_parent_ground_space_cd_identity} and
the normalization
\begin{align}
  \sum_{i_1}A^j_{i_1}A^{j\dagger}_{i_1}
    &=
    \Id,
  \label{eq:cpgsv21_block_diagonal_parent_ground_space_intersection_normalization}
\end{align}
one obtains
\begin{align}
  A^j_{i_{m+1}}E^jA^j_{i_1}
    &=
    \sum_k
    A^j_{i_{m+1}}C^j_kA^{j\dagger}_kA^j_{i_1},
  \label{eq:cpgsv21_block_diagonal_parent_ground_space_e_expand}\\
  &=
    D^j_{i_{m+1}}
    \left(\sum_k A^j_kA^{j\dagger}_k\right)A^j_{i_1},
  \label{eq:cpgsv21_block_diagonal_parent_ground_space_e_substitute}\\
  &=
    D^j_{i_{m+1}}A^j_{i_1},
  \label{eq:cpgsv21_block_diagonal_parent_ground_space_e_normalize}\\
  &=
    A^j_{i_{m+1}}C^j_{i_1}.
  \label{eq:cpgsv21_block_diagonal_parent_ground_space_e_conclude}
\end{align}
This calculation places $\psi$ in
$\bigoplus_j\mathcal G_{m+1}^{A_j}$ and proves the nontrivial inclusion in
Eq.~\ref{eq:cpgsv21_block_diagonal_parent_ground_space_intersection_identity}.

\section{The boundary condition and the global change of cut}

The formal one-step theorem assumes the simultaneous product-span identity
\begin{align}
  \operatorname{span}
  \{(A^1_w,\ldots,A^b_w):|w|=m\}
    &=
    \prod_{j=1}^bM_{D_j}(\mathbb C).
  \label{eq:cpgsv21_block_diagonal_parent_ground_space_product_span}
\end{align}
In the normalized BNT case, it also gives the internal direct sums
\begin{align}
  \bigvee_j\mathcal G_m^{A_j}
    &=
    \bigoplus_j\mathcal G_m^{A_j},
  \label{eq:cpgsv21_block_diagonal_parent_ground_space_internal_sum_m}\\
  \bigvee_j\mathcal G_{m+1}^{A_j}
    &=
    \bigoplus_j\mathcal G_{m+1}^{A_j}.
  \label{eq:cpgsv21_block_diagonal_parent_ground_space_internal_sum_succ}
\end{align}
Consequently, for every sufficiently large $m$ satisfying
Eq.~\ref{eq:cpgsv21_block_diagonal_parent_ground_space_product_span}, the
one-step conclusion is
\begin{align}
  \mathbb C^d\otimes
  \left(\bigoplus_j\mathcal G_m^{A_j}\right)
  \cap
  \left(\bigoplus_j\mathcal G_m^{A_j}\right)\otimes\mathbb C^d
    &=
    \bigoplus_j\mathcal G_{m+1}^{A_j}.
  \label{eq:cpgsv21_block_diagonal_parent_ground_space_formal_intersection}
\end{align}

The tuple-span, internal-sum, and one-step intersection theorems hold under the
source canonical data. For each block $j$, let $\Lambda^j$ be positive
definite and assume
\begin{align}
  \sum_aA^j_aA^{j\dagger}_a
    &=
    \Id,
  \label{eq:cpgsv21_block_diagonal_parent_ground_space_block_unital}\\
  \sum_aA^{j\dagger}_a\Lambda^jA^j_a
    &=
    \Lambda^j.
  \label{eq:cpgsv21_block_diagonal_parent_ground_space_block_dual_fixed}
\end{align}
These results require neither diagonality of $\Lambda^j$ nor the additional
identity $\sum_aA^{j\dagger}_aA^j_a=\Id$ for the same representative. Their
proof uses the positive dual fixed point only to construct a temporary
left-canonical gauge for the sharp base separation theorem. It transports the
tuple span back to the source gauge and propagates the span using
Eq.~\ref{eq:cpgsv21_block_diagonal_parent_ground_space_block_unital}.\footnote{Lean
declarations:
\leanid{MPSTensor.wordTupleSpanTop_threeBlock_mul_pred_of_blocksNotGaugePhaseEquiv_c1_of_dualFixedPoint},
\leanid{MPSTensor.wordTupleSpanTop_of_ge_of_bnt_directSum_unital_c1_pgvwc07_of_dualFixedPoint},
\leanid{MPSTensor.groundSpace_iSupIndep_of_ge_of_bnt_directSum_unital_c1_pgvwc07_of_dualFixedPoint},
and
\leanid{MPSTensor.pgvwc07_iSup_restriction_intersection_of_ge_of_bnt_directSum_unital_c1_pgvwc07_of_dualFixedPoint}.}
The earlier doubly normalized declarations remain available as
special-purpose siblings.

The normalization relaxation extends through open-boundary containment, unique
block decomposition, block-diagonal boundary representation, global change of
cut, periodic chain-space equality, kernel containment, and exact kernel
equality. Under the source-unital and positive dual fixed-point hypotheses, it
gives
\begin{align}
  \mathcal K_{N,L}(\widehat A)
    &=
    \bigvee_{j=1}^b\mathcal K_{N,L}(A_j),
  \label{eq:cpgsv21_block_diagonal_parent_ground_space_chain_decomposition}\\
  \ker H_L^{(N)}(\widehat A)
    &=
    \operatorname{span}
    \{V^{(N)}(A_j):1\leq j\leq b\}.
  \label{eq:cpgsv21_block_diagonal_parent_ground_space_kernel_span}
\end{align}
No source left-canonical assumption is needed.\footnote{Lean declarations:
\leanid{MPSTensor.chainGroundSpace_toTensorFromBlocks_le_iSup_groundSpace_of_ge_of_bnt_directSum_unital_c1_pgvwc07_of_dualFixedPoint},
\leanid{MPSTensor.exists_unique_sum_groundSpace_of_chainGroundSpace_toTensorFromBlocks_of_bnt_unital_c1_pgvwc07_of_dualFixedPoint},
\leanid{MPSTensor.exists_blockDiagonal_boundary_chainGroundSpace_of_global_cut_bnt_c1_pgvwc07_of_dualFixedPoint},
and
\leanid{MPSTensor.ker_parentHamiltonian_toTensorFromBlocks_eq_bntMPSVectorSpan_of_global_cut_bnt_c1_pgvwc07_of_dualFixedPoint}.}

Earlier conditional reductions isolated the source comparison after opening the
periodic boundary at a cut. Suppose
\begin{align}
  \psi
    &=
    \Gamma_N^{\widehat A}
    \left(\bigoplus_{j=1}^bX_j\right),
  \label{eq:cpgsv21_block_diagonal_parent_ground_space_boundary_representation}\\
  \psi
    &\in
    \mathcal K_{N,L}(\widehat A).
  \label{eq:cpgsv21_block_diagonal_parent_ground_space_boundary_state}
\end{align}
For a cyclic interval beginning at $i$ and crossing the cut, define
\begin{align}
  a
    &=
    i+L-N.
  \label{eq:cpgsv21_block_diagonal_parent_ground_space_wrapped_length}
\end{align}
The word $\beta$ records the part of the interval before the chosen cut, and
the word $\rho$ records the outside segment of length $N-L$. These blocked
words arise by opening the source boundary-index comparison.

The conditional theorem assumes that for every block $j$, every such outside
word $\rho$, and every word $\beta$ of length $a$, there is a matrix
$E_{j,i,\rho}$ such that
\begin{align}
  \mu_j^NX_jA^j_\beta A^j_\rho
    &=
    A^j_\beta E_{j,i,\rho}.
  \label{eq:cpgsv21_block_diagonal_parent_ground_space_complementary_identity}
\end{align}
Under the usual block-injectivity and normalization hypotheses, this condition
implies
\begin{align}
  \Gamma_N^{A_j}(\mu_j^NX_j)
    &\in
    \mathcal K_{N,L}(A_j)
    \qquad (1\leq j\leq b).
  \label{eq:cpgsv21_block_diagonal_parent_ground_space_component_membership}
\end{align}
Thus the trace representation of $\psi$ is not the obstruction. The periodic
component theorem derives
Eq.~\ref{eq:cpgsv21_block_diagonal_parent_ground_space_component_membership}
from
Eq.~\ref{eq:cpgsv21_block_diagonal_parent_ground_space_complementary_identity}.
The global change-of-cut proof avoids imposing the latter identity separately:
it translates the entire periodic state past the final $L_0$ sites and compares
the two open-boundary decompositions on the long complementary
segment.\footnote{Formal declarations:
\leanid{MPSTensor.block_boundary_intertwines_of_cyclicTranslate_sum_groundSpaceMap_eq},
\leanid{MPSTensor.exists_blockDiagonal_boundary_chainGroundSpace_of_global_cut_bnt_c1_pgvwc07_of_dualFixedPoint},
and
\leanid{MPSTensor.chainGroundSpace_toTensorFromBlocks_eq_iSup_of_global_cut_bnt_c1_pgvwc07_of_dualFixedPoint}.
The earlier conditional declarations are
\leanid{MPSTensor.blockDiagonal_boundary_component_chainGroundSpace_of_complementary_word_identities},
\leanid{MPSTensor.blockDiagonal_boundary_component_chainGroundSpace_of_complementary_word_identities_of_injective},
\leanid{MPSTensor.blockDiagonal_boundary_component_chainGroundSpace_of_pgvwc_comparison_of_injective},
\leanid{MPSTensor.exists_blockDiagonal_boundary_chainGroundSpace_of_complementary_identities_bnt_c1},
and
\leanid{MPSTensor.chainGroundSpace_toTensorFromBlocks_eq_iSup_and_iSupIndep_of_complementary_identities}.}

The live blueprint records the block-diagonal periodic-boundary ground space,
the forward inclusion, the block-diagonal intersection identity from
Theorem~12 of Ref.~\cite{PerezGarcia2007MPSReps}, the parent-ground-space
containment, and the final span statement. It distinguishes the opened-boundary
coordinates from the source boundary indices and records the unconditional
global-cut conclusion at the source length bound.
% Blueprint labels: def:parent_hamiltonian_ground_space_bnt,
% lem:isup_chain_ground_space_block_le_assembled,
% thm:block_diagonal_ground_space_intersection,
% thm:parent_ground_space_le_bnt_span, thm:gs_eq_bnt_span.

The sum of the blockwise chain spaces means
\begin{align}
  \sum_{j=1}^b\mathcal K_{N,L}(A_j)
    &=
    \left\{\sum_{j=1}^b\psi_j:
      \psi_j\in\mathcal K_{N,L}(A_j)\right\}.
  \label{eq:cpgsv21_block_diagonal_parent_ground_space_space_sum}
\end{align}
The proved source-range statement is
\begin{align}
  &\left[b\geq2
    \ \wedge\
    L_0\geq1
    \ \wedge\
    (\forall j,\ \mu_j\neq0)
    \ \wedge\
    L_0<L
    \ \wedge\
    L\geq3(b-1)(L_0+1)+1
    \ \wedge\
    N\geq L+L_0\right]
  \nonumber\\
  &\qquad\Longrightarrow
  \mathcal K_{N,L}(\widehat A)
  =
  \sum_{j=1}^b\mathcal K_{N,L}(A_j).
  \label{eq:cpgsv21_block_diagonal_parent_ground_space_source_range_decomposition}
\end{align}

An arbitrary trace-form vector need not belong to the periodic chain space:
\begin{align}
  \mathcal G_N^A
    &\nsubseteq
    \mathcal K_{N,L}(A)
    \quad\text{in general}.
  \label{eq:cpgsv21_block_diagonal_parent_ground_space_false_trace_inclusion}
\end{align}
If a cyclic window crosses the chosen virtual boundary, a vector
$\Gamma_N^A(X)$ has coefficients of the form
\begin{align}
  \tr\!\left(A^{\omega_{\mathrm{tail}}}
    XA^{\omega_{\mathrm{head}}}A^\rho\right),
  \label{eq:cpgsv21_block_diagonal_parent_ground_space_crossing_trace}
\end{align}
where $\omega_{\mathrm{tail}}\omega_{\mathrm{head}}$ is the word observed in
the cyclic window and $\rho$ is the outside word. For general $X$, this
expression cannot be written as
\begin{align}
  \tr\!\left(A^{\omega_{\mathrm{tail}}}
    A^{\omega_{\mathrm{head}}}Y_\rho\right)
  \label{eq:cpgsv21_block_diagonal_parent_ground_space_local_trace}
\end{align}
with one matrix $Y_\rho$ independent of the window word. The missing step is
therefore a boundary-condition comparison, not the false inclusion in
Eq.~\ref{eq:cpgsv21_block_diagonal_parent_ground_space_false_trace_inclusion}.

The source theorem uses the multi-block hypothesis, a positive common
injectivity length, nonzero weights, and the length range of Theorem~12 in
Ref.~\cite{PerezGarcia2007MPSReps}:
\begin{align}
  b
    &\geq
    2,
  \label{eq:cpgsv21_block_diagonal_parent_ground_space_hyp_blocks}\\
  L_0
    &\geq
    1,
  \label{eq:cpgsv21_block_diagonal_parent_ground_space_hyp_injectivity}\\
  \mu_j
    &\neq
    0
    \qquad\text{for every }j,
  \label{eq:cpgsv21_block_diagonal_parent_ground_space_hyp_weights}\\
  L_0
    &<
    L,
  \label{eq:cpgsv21_block_diagonal_parent_ground_space_hyp_local_strict}\\
  L
    &\geq
    3(b-1)(L_0+1)+1,
  \label{eq:cpgsv21_block_diagonal_parent_ground_space_hyp_local_bound}\\
  N
    &\geq
    L+L_0.
  \label{eq:cpgsv21_block_diagonal_parent_ground_space_hyp_global_bound}
\end{align}
For $b\geq2$ and $L_0\geq1$, the source bound implies
\begin{align}
  L
    &\geq
    3L_0+4
    >
    L_0.
  \label{eq:cpgsv21_block_diagonal_parent_ground_space_implied_strict_bound}
\end{align}
The hypothesis $L_0<L$ remains explicit because the proof uses one-block
uniqueness in that range. The one-block case is governed by the injective
parent-Hamiltonian theorem, not by the degenerate block-injective boundary
theorem.

The former formal proof used the more conservative range
\begin{align}
  L
    &\geq
    (L_0+1)+3(b-1)(L_0+1)+1.
  \label{eq:cpgsv21_block_diagonal_parent_ground_space_old_length_bound}
\end{align}
The extra $(L_0+1)$ is no longer needed. Once simultaneous tuples span at
\begin{align}
  q
    &=
    3(b-1)(L_0+1),
  \label{eq:cpgsv21_block_diagonal_parent_ground_space_base_span_length}
\end{align}
the unital normalization propagates the homogeneous tuple span from $q$ to
every greater length. Given a target tuple $(M_j)_j$, represent
$((A^j_a)^\dagger M_j)_j$ at length $q$, prefix $A^j_a$, and sum over $a$.

The periodic-boundary reductions have three forms: a decomposition into
blockwise vectors with periodic boundary conditions; a block-diagonal boundary
representation whose components already lie in the corresponding blockwise
periodic spaces; and the boundary-crossing coordinate formulation in
Eq.~\ref{eq:cpgsv21_block_diagonal_parent_ground_space_complementary_identity}.\footnote{Lean
declarations:
\leanid{MPSTensor.chainGroundSpace_toTensorFromBlocks_le_iSup_of_boundary_decomposition},
\leanid{MPSTensor.chainGroundSpace_toTensorFromBlocks_le_iSup_of_blockDiagonal_boundary_groundSpaceMap},
\leanid{MPSTensor.chainGroundSpace_toTensorFromBlocks_eq_iSup_and_iSupIndep_of_bnt_c1_boundary_decomposition},
\leanid{MPSTensor.chainGroundSpace_toTensorFromBlocks_eq_iSup_and_iSupIndep_of_bnt_c1_blockBoundary},
and
\leanid{MPSTensor.chainGroundSpace_toTensorFromBlocks_eq_iSup_and_iSupIndep_of_complementary_identities}.}
The global change-of-cut theorem discharges the boundary-crossing hypothesis at
the source length constant.

For a block-diagonal boundary matrix
\begin{align}
  X
    &=
    \bigoplus_jX_j,
  \label{eq:cpgsv21_block_diagonal_parent_ground_space_block_boundary}
\end{align}
the boundary formula is
\begin{align}
  \Gamma_N^{\widehat A}(X)(i_1,\ldots,i_N)
    &=
    \sum_j\mu_j^N
    \tr(A^j_{i_1}\cdots A^j_{i_N}X_j),
  \label{eq:cpgsv21_block_diagonal_parent_ground_space_weighted_trace_sum}\\
  &=
    \sum_j\mu_j^N
    \Gamma_N^{A_j}(X_j)(i_1,\ldots,i_N).
  \label{eq:cpgsv21_block_diagonal_parent_ground_space_weighted_gamma_sum}
\end{align}
For the remaining argument, abbreviate
\begin{align}
  \Gamma_j(X)
    &=
    \Gamma_N^{A_j}(X),
  \label{eq:cpgsv21_block_diagonal_parent_ground_space_gamma_abbreviation}\\
  V_j
    &=
    V^{(N)}(A_j).
  \label{eq:cpgsv21_block_diagonal_parent_ground_space_vector_abbreviation}
\end{align}

The periodic-boundary proof first establishes
\begin{align}
  \psi\in\mathcal K_{N,L}(\widehat A)
    &\Longrightarrow
    \exists(X_j)_{j=1}^b,\
    \psi=\sum_{j=1}^b\Gamma_j(X_j),
  \label{eq:cpgsv21_block_diagonal_parent_ground_space_first_implication_decompose}\\
  \psi\in\mathcal K_{N,L}(\widehat A)
    &\Longrightarrow
    \Gamma_j(X_j)\in\mathcal K_{N,L}(A_j)
    \quad (1\leq j\leq b).
  \label{eq:cpgsv21_block_diagonal_parent_ground_space_first_implication_components}
\end{align}
For arbitrary $X_1,\ldots,X_b$, one-block uniqueness then gives
\begin{align}
  \left[\Gamma_j(X_j)\in\mathcal K_{N,L}(A_j)
    \text{ for every }j\right]
    &\Longrightarrow
    \exists(c_j)_{j=1}^b,\
    \Gamma_j(X_j)=c_jV_j
    \text{ for every }j.
  \label{eq:cpgsv21_block_diagonal_parent_ground_space_second_implication}
\end{align}
Combining these implications yields
\begin{align}
  \psi
    &=
    \sum_{j=1}^b\Gamma_j(X_j),
  \label{eq:cpgsv21_block_diagonal_parent_ground_space_span_composition_first}\\
  &=
    \sum_{j=1}^bc_jV_j,
  \label{eq:cpgsv21_block_diagonal_parent_ground_space_span_composition_second}\\
  \psi
    &\in
    \operatorname{span}\{V_j:j=1,\ldots,b\}.
  \label{eq:cpgsv21_block_diagonal_parent_ground_space_span_conclusion}
\end{align}

\section{Conclusion}

This note does not identify a source error. The formal development proves the
block-diagonal periodic-boundary comparison without a short-tail span
hypothesis or an external boundary-comparison hypothesis in the source range
\begin{align}
  b
    &\geq
    2,
  \label{eq:cpgsv21_block_diagonal_parent_ground_space_conclusion_blocks}\\
  L_0
    &\geq
    1,
  \label{eq:cpgsv21_block_diagonal_parent_ground_space_conclusion_injectivity}\\
  L
    &\geq
    3(b-1)(L_0+1)+1,
  \label{eq:cpgsv21_block_diagonal_parent_ground_space_conclusion_local_length}\\
  N
    &\geq
    L+L_0.
  \label{eq:cpgsv21_block_diagonal_parent_ground_space_conclusion_global_length}
\end{align}
For each block $j$, it proves
\begin{align}
  \psi\in\mathcal K_{N,L}(\widehat A)
    &\Longrightarrow
    \Gamma_N^{A_j}(\mu_j^NX_j)
    \in
    \mathcal K_{N,L}(A_j).
  \label{eq:cpgsv21_block_diagonal_parent_ground_space_conclusion_component}
\end{align}

The earlier conditional reduction recorded by pull request~\#2724 and the
boundary-trace implication remain useful intermediate results. The latter has
the following form. Assume the common word-span property at middle-word length
$m$. For every boundary-crossing interval, let $\beta$ be the word before the
cut, let $\rho$ be the outside word, and let $w$ be a middle word of length
$m$. If the two boundary expressions obtained from Theorem~12 of
Ref.~\cite{PerezGarcia2007MPSReps} satisfy
\begin{align}
  \sum_j
  \tr\!\left(A^j_\beta C^j_\rho A^j_w\right)
    &=
    \sum_j
    \tr\!\left((\mu_j^NX_jA^j_\beta)A^j_\rho A^j_w\right),
  \label{eq:cpgsv21_block_diagonal_parent_ground_space_trace_decomposition}
\end{align}
then every block in the block-injective range satisfies
\begin{align}
  \Gamma_N^{A_j}(\mu_j^NX_j)
    &\in
    \mathcal K_{N,L}(A_j).
  \label{eq:cpgsv21_block_diagonal_parent_ground_space_trace_component}
\end{align}
\footnote{Lean declaration:
\leanid{MPSTensor.blockDiagonal_boundary_component_chainGroundSpace_of_trace_decomposition_of_injective}.}

An explicit-boundary version fixes
\begin{align}
  \psi
    &=
    \Gamma_N^{\widehat A}
    \left(\bigoplus_jX_j\right).
  \label{eq:cpgsv21_block_diagonal_parent_ground_space_explicit_boundary}
\end{align}
With this boundary representation, the same equality of boundary trace
expressions implies
\begin{align}
  \Gamma_N^{A_j}(\mu_j^NX_j)
    &\in
    \mathcal K_{N,L}(A_j)
  \label{eq:cpgsv21_block_diagonal_parent_ground_space_explicit_component}
\end{align}
for every block $j$, without a span assumption at the short
boundary-crossing tail length $N-i$.\footnote{Lean declaration:
\leanid{MPSTensor.exists_blockDiagonal_boundary_chainGroundSpace_of_trace_decomposition_of_boundary}.}

At the fixed wrapped-tail length supplied by one boundary-crossing cyclic
window, the local constraint itself gives the boundary trace expressions once
the block-diagonal boundary representation has been fixed. More precisely, let
$N<i+L$, let $\rho$ be an outside word of length $N-L$, and let $\alpha$ be a
wrapped-tail word of length $N-i$. There are matrices $C^j_{i,\rho}$ such that
\begin{align}
  \sum_j
  \tr\!\left(A^j_\beta C^j_{i,\rho}A^j_\alpha\right)
    &=
    \sum_j
    \tr\!\left((\mu_j^NX_jA^j_\beta)A^j_\rho A^j_\alpha\right).
  \label{eq:cpgsv21_block_diagonal_parent_ground_space_fixed_tail_trace}
\end{align}
This fixed-tail statement is distinct from the middle-word trace hypothesis at
a common product-spanning length. The fixed cyclic window supplies the
wrapped-tail length $N-i$, whereas the separating trace-decomposition
implication uses a common product-span length $m$.\footnote{Formal declaration:
\leanid{MPSTensor.blockDiagonal_boundary_crossing_trace_decomposition_of_boundary}.}

In the normalized BNT source range, the common middle-word span is no longer an
external assumption. If
\begin{align}
  m
    &=
    3(b-1)(L_0+1),
  \label{eq:cpgsv21_block_diagonal_parent_ground_space_middle_length}\\
  L
    &\geq
    m+1,
  \label{eq:cpgsv21_block_diagonal_parent_ground_space_middle_local_bound}
\end{align}
then normalized block separation gives
\begin{align}
  \operatorname{span}
  \{(A^1_w,\ldots,A^b_w):|w|=m\}
    &=
    \prod_jM_{D_j}(\mathbb C).
  \label{eq:cpgsv21_block_diagonal_parent_ground_space_middle_product_span}
\end{align}
Consequently, after choosing a block-diagonal boundary representation, the
finite-range trace-decomposition implication retains only
Eq.~\ref{eq:cpgsv21_block_diagonal_parent_ground_space_trace_decomposition} as
an external input.\footnote{Lean declarations:
\leanid{MPSTensor.exists_blockDiagonal_boundary_chainGroundSpace_of_trace_decomposition_bnt_c1_span}
and
\leanid{MPSTensor.chainGroundSpace_toTensorFromBlocks_eq_iSup_and_iSupIndep_of_trace_decomposition_bnt_c1_span}.}

The comparison of $C^j$ and $D^j$ after opening the boundary is closer to the
proof of Theorem~12 in Ref.~\cite{PerezGarcia2007MPSReps}. The local words
$\beta$ and $\rho$ arise after opening and blocking the cyclic comparison.
They reindex the source end-site comparison $i_1,i_{m+1}$ and introduce no
additional source hypothesis or terminology.

For a fixed boundary-crossing interval beginning at site $i$, the local
block-sum constraint and the word-tuple span at wrapped-tail length $N-i$ give
matrices $C^j_{i,\rho}$ satisfying
\begin{align}
  A^j_\beta C^j_{i,\rho}
    &=
    ((\mu_j^NX_j)A^j_\beta)A^j_\rho.
  \label{eq:cpgsv21_block_diagonal_parent_ground_space_compatibility}
\end{align}
Once
Eq.~\ref{eq:cpgsv21_block_diagonal_parent_ground_space_compatibility} is
available, the normalized $E^j$ calculation is complete. Define
\begin{align}
  E^j_i
    &=
    \sum_\rho C^j_{i,\rho}(A^j_\rho)^\dagger.
  \label{eq:cpgsv21_block_diagonal_parent_ground_space_crossing_e}
\end{align}
The resulting identities are
\begin{align}
  (\mu_j^NX_j)A^j_\beta
    &=
    A^j_\beta E^j_i,
  \label{eq:cpgsv21_block_diagonal_parent_ground_space_crossing_intertwiner}\\
  A^j_\beta C^j_{i,\rho}
    &=
    A^j_\beta E^j_iA^j_\rho.
  \label{eq:cpgsv21_block_diagonal_parent_ground_space_crossing_factorization}
\end{align}
The local boundary formula is
\begin{align}
  E_{j,i,\rho}
    &=
    E^j_iA^j_\rho.
  \label{eq:cpgsv21_block_diagonal_parent_ground_space_local_boundary_formula}
\end{align}
These identities, the local formula, and its existential consequence are
formalized by the corresponding complementary-word declarations.\footnote{Lean
declarations:
\leanid{MPSTensor.pgvwc07_complementary_word_cde_identities_of_block_boundary_trace_decomposition},
\leanid{MPSTensor.pgvwc07_complementary_word_boundary_identities_formula_of_compatibility},
and
\leanid{MPSTensor.pgvwc07_complementary_word_boundary_identities_of_compatibility}.}

The crossing comparison first treats membership in the sum of block spaces,
then specializes to a fixed block-diagonal boundary representation, and
finally yields the required blockwise chain-space statement. The current BNT
specialization supplies that boundary representation, so the crossing
comparison is not an independent trace-comparison assumption.\footnote{Lean
declarations:
\leanid{MPSTensor.blockDiagonal_boundary_crossing_pgvwc_comparison_of_sum_mem_iSup},
\leanid{MPSTensor.blockDiagonal_boundary_crossing_pgvwc_comparison_of_chainGroundSpace},
\leanid{MPSTensor.exists_blockDiagonal_boundary_chainGroundSpace_of_crossing_pgvwc_comparison_of_boundary},
and
\leanid{MPSTensor.exists_blockDiagonal_boundary_chainGroundSpace_of_crossing_pgvwc_comparison_bnt_c1}.}

The earlier crossing-span route required both a crossing-tail word-span
hypothesis and a block-diagonal boundary representation for
$\psi\in\mathcal K_{N,L}(\widehat A)$. In the normalized BNT specialization,
the block-diagonal boundary theorem supplies the representation. The
crossing-tail span must not be confused with the large-length BNT product-span
bound: if the boundary-crossing interval begins at $i=N-1$, then
\begin{align}
  N-i
    &=
    1.
  \label{eq:cpgsv21_block_diagonal_parent_ground_space_short_crossing_tail}
\end{align}
Thus a lower bound on $L$ alone does not supply the required tail span.

The block-diagonal boundary representation is not an instance of the
boundary-crossing comparison. It follows from the source-shaped, span-based
one-step intersection theorem. Its block components belong to the
open-boundary spaces $\mathcal G_N^{A_j}$, and its proof does not use the
periodic-boundary comparison at crossing windows.\footnote{Lean declarations:
\leanid{MPSTensor.exists_blockDiagonal_boundary_of_chainGroundSpace_toTensorFromBlocks_of_bnt_unital_c1_pgvwc07_of_dualFixedPoint}
and
\leanid{MPSTensor.pgvwc07_iSup_restriction_intersection_of_ge_of_bnt_directSum_unital_c1_pgvwc07_of_dualFixedPoint}.}
Earlier conditional block-diagonal parent-space theorems retain the
periodic-boundary comparison as an explicit hypothesis, whereas the global
change-of-cut theorem discharges it. Homogeneous tuple-span propagation removes
the former extra $(L_0+1)$ from the normalized BNT product-span range.

The periodic-boundary target of
\href{https://github.com/LionSR/TNLean/issues/3597}{issue \#3597} is therefore
complete at the source range. The length-constant strengthening tracked by
\href{https://github.com/LionSR/TNLean/issues/4195}{issue \#4195} is also
complete. Its hypotheses are exactly
\begin{align}
  b
    &\geq
    2,
  \label{eq:cpgsv21_block_diagonal_parent_ground_space_final_blocks}\\
  L_0
    &\geq
    1,
  \label{eq:cpgsv21_block_diagonal_parent_ground_space_final_injectivity}\\
  L
    &\geq
    3(b-1)(L_0+1)+1,
  \label{eq:cpgsv21_block_diagonal_parent_ground_space_final_local_length}\\
  N
    &\geq
    L+L_0.
  \label{eq:cpgsv21_block_diagonal_parent_ground_space_final_global_length}
\end{align}
Under these hypotheses, the global-cut theorem gives
\begin{align}
  \mathcal K_{N,L}(\widehat A)
    &=
    \sum_{j=1}^b\mathcal K_{N,L}(A_j).
  \label{eq:cpgsv21_block_diagonal_parent_ground_space_final_chain_decomposition}
\end{align}
Applying the one-block uniqueness theorem to each block then gives the final
degenerate parent-Hamiltonian ground-space equality
\begin{align}
  \ker H_L^{(N)}(\widehat A)
    &=
    \operatorname{span}
    \{V^{(N)}(A_j):1\leq j\leq b\}.
  \label{eq:cpgsv21_block_diagonal_parent_ground_space_final_kernel}
\end{align}

\bibliographystyle{alpha}
\bibliography{references}

\end{document}
